\documentclass[9pt,a4paper]{article}

\usepackage[utf8]{inputenc}
\usepackage[T1]{fontenc}
\usepackage{mathpazo}
\usepackage[scaled=0.92]{helvet}
\usepackage[margin=0.5in,top=0.4in,bottom=0.4in]{geometry}
\usepackage{amsmath,amssymb}
\usepackage{booktabs}
\usepackage{tabularx}
\usepackage{multicol}
\usepackage{enumitem}
\usepackage{setspace}

\setstretch{1.0}
\setlength{\parindent}{0pt}
\setlength{\parskip}{0.15em}
\setlength{\columnsep}{0.4cm}
\setlength{\columnseprule}{0.3pt}
\setlength{\emergencystretch}{2em}

\setlist[itemize]{leftmargin=1em,noitemsep,topsep=0.5pt,parsep=0pt,itemsep=0.5pt}
\setlist[enumerate]{leftmargin=1.1em,noitemsep,topsep=0.5pt,parsep=0pt,itemsep=0.5pt}
\renewcommand{\arraystretch}{1.1}

\newcommand{\mysect}[1]{\par\vspace{4pt}\noindent{\small\bfseries\sffamily #1}\par\vspace{-3pt}\noindent\rule{\linewidth}{0.4pt}\par\vspace{2pt}\noindent\ignorespaces}
\newcommand{\h}[1]{\textbf{#1}}

\pagestyle{empty}

\begin{document}

\begin{center}
	{\large\bfseries\sffamily Cheat Sheet: Fungsi}\\[1pt]
	{\scriptsize Matematika Diskrit \textbullet\ TRPL Politeknik Negeri Batam}
\end{center}
\vspace{-8pt}

\begin{multicols}{2}

	\mysect{1. Definisi Fungsi}
	$f: A\to B$: setiap $a\in A$ dipasangkan dengan \h{tepat satu} $b\in B$. Notasi $f(a)=b$.
	\begin{itemize}
		\item \h{Domain} $A$: daerah asal.
		\item \h{Kodomain} $B$: tempat hasil boleh berada.
		\item \h{Range}: hasil yang benar-benar muncul. $\mathrm{ran}(f)\subseteq B$.
	\end{itemize}

	\mysect{2. Cek Fungsi: 3 Pertanyaan}
	\begin{enumerate}
		\item Apa saja masukannya?
		\item Tiap masukan punya berapa hasil? Lebih dari satu $\Rightarrow$ bukan fungsi.
		\item Ada masukan tanpa hasil? $\Rightarrow$ fungsi parsial.
	\end{enumerate}
	Kasus kampus: $M3\to MTK$, $M3\to RPL$ \textbf{bukan} fungsi (dua hasil). Jika M3 belum punya dosen wali, relasinya parsial.

	\mysect{3. Sifat Fungsi}
	{\footnotesize
		\begin{tabular}{@{}l l@{}}
			\toprule
			Sifat     & Arti                                       \\
			\midrule
			Injektif  & $a_1\neq a_2\Rightarrow f(a_1)\neq f(a_2)$ \\
			Surjektif & $\forall b\in B,\exists a\in A:f(a)=b$     \\
			Bijektif  & Injektif \emph{dan} surjektif              \\
			\bottomrule
		\end{tabular}}

	Cek cepat: injektif butuh $|A|\le |B|$; surjektif butuh $|A|\ge |B|$. Syarat perlu, belum cukup.

	\mysect{4. Invers Fungsi}
	$f(a)=b\Rightarrow f^{-1}(b)=a$. \h{Syarat:} $f$ bijektif.\\
	\h{Contoh:} $f(x)=2x+3\Rightarrow f^{-1}(x)=(x-3)/2$.\\
	\h{Langkah:} (1) cek bijektif, (2) $y=f(x)$, (3) cari $x$, (4) tukar $y\to x$, (5) uji satu angka.\\
	\h{Peringatan:} $f(x)=x^2$ tanpa batas domain tidak punya invers ($9\to 3$ atau $-3$?). Batasi $x\ge 0$.

	\mysect{5. Komposisi Fungsi}
	$(g\circ f)(x)=g(f(x))$: $f$ dulu, hasilnya masuk ke $g$.\\
	\h{Syarat:} $\mathrm{ran}(f)\subseteq\mathrm{dom}(g)$.\\
	\h{Contoh:} $f(x)=3x+1$, $g(x)=x^2$. $g\circ f(2)=g(7)=49$; $f\circ g(2)=f(4)=13$. Urutan penting.\\
	\h{Domain:} $f(x)=x-5$, $g(x)=\sqrt{x}\Rightarrow g\circ f(x)=\sqrt{x-5}$, hanya $x\ge 5$.

	\mysect{6. Fungsi Khusus}
	{\footnotesize
		\begin{tabular}{@{}l l l@{}}
			\toprule
			Fungsi                    & Arti                     & Contoh                                               \\
			\midrule
			Lantai $\lfloor x\rfloor$ & bulat ke bawah           & $\lfloor 3{,}7\rfloor=3$, $\lfloor -1{,}2\rfloor=-2$ \\
			Langit $\lceil x\rceil$   & bulat ke atas            & $\lceil 3{,}2\rceil=4$                               \\
			Modulo $a\bmod n$         & sisa bagi                & $17\bmod 5=2$                                        \\
			Faktorial $n!$            & $n(n-1)\cdots 1$, $0!=1$ & $5!=120$                                             \\
			\bottomrule
		\end{tabular}}

	Hubungan: $a=n\cdot q+r$, $q=\lfloor a/n\rfloor$, $r=a\bmod n$.\\
	Jumlah halaman $=\lceil total/10\rceil$: $47\to 5$, $100\to 10$, $101\to 11$.\\
	\h{Hati-hati:} $-7\bmod 3=2$ (matematika), tapi $-1$ di C/Java.

	\mysect{7. Fungsi Rekursif}
	Panggil diri sendiri. \h{Dua bagian:} basis (berhenti) + rekurens (memecah). Tanpa basis $\Rightarrow$ loop tak berhenti.\\
	\h{Faktorial:} $n!=n\cdot (n-1)!$, $0!=1$. Pseudocode: $fakt(n)=1$ jika $n=0$; $fakt(n)=n\cdot fakt(n-1)$ jika $n>0$.\\
	\h{Deret:} $S(n)=n+S(n-1)$, $S(0)=0$.\\
	\h{Fibonacci:} $F(n)=F(n-1)+F(n-2)$, $F(0)=0$, $F(1)=1$: $0,1,1,2,3,5,8,\dots$\\
	Eksponensial $\Rightarrow$ pakai loop atau memoization.

	\mysect{8. Studi Kasus: Sandi Geser}
	A$=0$, $\dots$, Z$=25$. Rumus $g(x)=(x+3)\bmod 26$.\\
	Bijektif $\Rightarrow$ invers $g^{-1}(x)=(x-3)\bmod 26=(x+23)\bmod 26$.\\
	\h{Contoh:} RPL $=(17,15,11)\to(20,18,14)=$ USO. Balik: USO $\to$ RPL.\\
	Komposisi: $(g\circ g)(x)=(x+6)\bmod 26$.

	\mysect{9. Kasus Uji}
	{\footnotesize
		\begin{tabular}{@{}l l l@{}}
			\toprule
			Jenis       & Masukan    & Hasil \\
			\midrule
			Normal      & $(47,10)$  & 5     \\
			Batas habis & $(100,10)$ & 10    \\
			Batas lebih & $(101,10)$ & 11    \\
			Batas nol   & $(0,10)$   & 0     \\
			Luar domain & $(50,0)$   & galat \\
			\bottomrule
		\end{tabular}}

	\mysect{10. Kesalahan Umum}
	{\footnotesize
		\begin{tabular}{@{}l l@{}}
			\toprule
			Salah                 & Benar                           \\
			\midrule
			Semua relasi = fungsi & Fungsi: 1 masukan $\to$ 1 hasil \\
			Range $=$ kodomain    & Range $\subseteq$ kodomain      \\
			$f^{-1}(x)=1/f(x)$    & $f^{-1}$ membalik pasangan      \\
			$g\circ f=f\circ g$   & Urutan penting, $f$ dulu        \\
			Rekursi tanpa basis   & Tidak akan berhenti             \\
			\bottomrule
		\end{tabular}}

	\mysect{11. Lembar Rujukan Cepat}
	{\footnotesize
		\begin{tabular}{@{}l l@{}}
			\toprule
			Butuh           & Rumus                                      \\
			\midrule
			Cek fungsi      & tiap $a$: tepat 1 hasil                    \\
			Injektif        & $a_1\neq a_2\Rightarrow f(a_1)\neq f(a_2)$ \\
			Surjektif       & $\mathrm{ran}(f)=B$                        \\
			Bijektif        & injektif + surjektif                       \\
			Invers          & $y=f(x)$, tukar, cari $x$                  \\
			Komposisi       & $g(f(x))$, $f$ dulu                        \\
			Lantai / Langit & $\lfloor x\rfloor$ / $\lceil x\rceil$      \\
			Modulo          & $a\bmod n = a-n\lfloor a/n\rfloor$         \\
			Rekursi         & basis + rekurens                           \\
			\bottomrule
		\end{tabular}}

\end{multicols}

\end{document}
